\appendices
\section{理论与离散}
\label{chapter:S1}

% block-id: S1.round_derivation.heading
\hypertarget{S1.round_derivation.heading}{}
\subsection{圆线径向方程与归一化}

% block-id: S1.round_derivation.p1
\hypertarget{S1.round_derivation.p1}{}
长直圆线内部的轴向电流满足 (1/r)d(r dJz/dr)/dr=κ²Jz。轴心有限性确定径向函数为 I\(_0\)(κr)，总复电流约束与积分关系 \(\int\)\(_0\)\(^a\)rI\(_0\)(κr)dr=aI\(_1\)(κa)/κ 确定其系数。表面轴向电场与总电流之比给出单位长度内部阻抗；对电流密度独立执行焦耳积分，可核对阻抗实部与损耗之间的关系。推导采用与正文一致的 RMS 相量。

% block-id: S1.sharing_derivation.heading
\hypertarget{S1.sharing_derivation.heading}{}
\subsection{复电流均流条件的推导}

% block-id: S1.sharing_derivation.p1
\hypertarget{S1.sharing_derivation.p1}{}
若等流解存在，将 i=(It/N)1 代入端口方程，即得向量 (It/N)Z1+e 的全部分量相同，投影条件因而必要。反之，若投影条件成立，该向量可写为某一标量 V 与 1 的乘积，等流向量同时满足端口方程和总电流约束；在增广系统解唯一时，它即为实际解，条件因而充分。数值结果以 sqrt(mean(|ik−It/N|²))/|It/N| 表示归一化复电流 RMS 偏差，并同时展示逐股幅值与相位。

% block-id: S1.dtm_derivation.heading
\hypertarget{S1.dtm_derivation.heading}{}
\subsection{股线身份与截面位置的变换}

% block-id: S1.dtm_derivation.p1
\hypertarget{S1.dtm_derivation.p1}{}
令 Ps 将固定物理股线身份映射到第 s 个截面的空间位置，则局部电流为 Ps i，转换回身份坐标的局部电压为 Ps\(^{\mathrm T}\)ZsPs i。按段长 Δzs 串联相加得到 \eqref{eq:E06}。在这一表示中，位置与身份分别保存，K、2K、4K 的分段比较使用同一轨迹和同一组股线。周期、间隙与曲率用于检查轨迹，实际弧长和铜体积用于比较材料用量。

% block-id: S1.equation.E06
\hypertarget{S1.equation.E06}{}
\hypertarget{E06}{}\begin{equation}Z_{\mathrm{tot}}=\sum_{s=1}^{K}\Delta z_s P_s^T Z_s P_s\tag{7}\label{eq:E06}\end{equation}

% block-id: S1.peec_derivation.heading
\hypertarget{S1.peec_derivation.heading}{}
\subsection{纵向细丝基函数与核积分}

% block-id: S1.peec_derivation.p1
\hypertarget{S1.peec_derivation.p1}{}
第 α 条体积细丝中的电流密度取为常量 iα/Aα，方向沿轴向。由矢势体积积分得到测试细丝内的感应电场，再沿长度积分，可得对称部分电感矩阵。自项在重合点具有可积奇性，使用专门的自项精度控制；互项采用独立积分阶数。欧姆电阻与 jωL 组成耦合矩阵，并与公共端电压和总电流约束联合求解。所用基函数决定了这一模型能够表示的电流方向。

% block-id: S1.triangle_derivation.heading
\hypertarget{S1.triangle_derivation.heading}{}
\subsection{线性三角形的电流与功率积分}

% block-id: S1.triangle_derivation.p1
\hypertarget{S1.triangle_derivation.p1}{}
线性三角形内以 J=\(\sum\)\(_a\)NaJa 插值复电流密度。每个形函数的面积积分为 AT/3，因而可直接得到单元复电流；焦耳损耗使用 \eqref{eq:E09} 的一致质量矩阵计算。材料标记确定单元所属股线、回流或空气，坐标转换为米后求面积。按材料求和所得电流和损耗与原始场路标量独立核对。

% block-id: S1.equation.E09
\hypertarget{S1.equation.E09}{}
\hypertarget{E09}{}
\begin{equation}
\begin{gathered}
M_T=\frac{A_T}{12}\begin{bmatrix}2&1&1\\1&2&1\\1&1&2\end{bmatrix} \\
P_T=\mathbf J_T^H M_T\mathbf J_T/\sigma
\end{gathered}
\tag{8}\label{eq:E09}
\end{equation}

% block-id: S1.method.sources
\hypertarget{S1.method.sources}{}
圆线推导采用修正 Bessel 函数的定义与积分关系；纵向耦合采用部分元件等效电路的电阻和部分电感表述。\cite{BESSEL,RUEHLI1974}
